% !TeX program = xelatex
\documentclass[UTF8,a4paper,zihao=-4,fontset=fandol]{ctexart}

\usepackage[left=28mm,right=26mm,top=27mm,bottom=25mm,headheight=16pt]{geometry}
\usepackage{setspace}
\usepackage{enumitem}
\usepackage{booktabs}
\usepackage{array}
\usepackage{tabularx}
\usepackage{caption}
\usepackage{fancyhdr}
\usepackage{xurl}
\usepackage[hidelinks,bookmarksnumbered]{hyperref}

\newcommand{\cmsversion}{0.4.5}
\newcommand{\menu}[2]{\textbf{#1}\,\ensuremath{\rightarrow}\,\textbf{#2}}
\newcommand{\code}[1]{\texttt{\detokenize{#1}}}
\newcolumntype{Y}{>{\raggedright\arraybackslash}X}
\newenvironment{manualtable}[1]
  {\par\medskip\noindent\begin{minipage}{\linewidth}\small\captionof{table}{#1}\centering}
  {\end{minipage}\par\medskip}

\setstretch{1.42}
\setlength{\parindent}{2\ccwd}
\setlength{\parskip}{0pt}
\setlength{\emergencystretch}{2em}
\setlength{\tabcolsep}{5pt}
\renewcommand{\arraystretch}{1.24}
\setlist[enumerate]{label=\arabic*.,leftmargin=2.7em,topsep=3pt,itemsep=2pt,parsep=0pt}
\setlist[itemize]{leftmargin=2.7em,topsep=3pt,itemsep=2pt,parsep=0pt}
\captionsetup{font=small,labelfont=bf,labelsep=quad,skip=5pt,hypcap=false}
\ctexset{
  section={format=\zihao{3}\heiti\raggedright,beforeskip=1.4ex plus .3ex,afterskip=.8ex},
  subsection={format=\zihao{4}\heiti\raggedright,beforeskip=1.1ex plus .3ex,afterskip=.5ex},
}
\setcounter{tocdepth}{1}
\renewcommand{\contentsname}{目\quad 录}
\hypersetup{
  pdftitle={DSSAST CMS 管理员操作手册},
  pdfsubject={赛事、学习笔记、新闻公告与秀米同步的 WordPress 操作指南},
  pdfdisplaydoctitle=true
}

\fancypagestyle{manual}{
  \fancyhf{}
  \fancyhead[L]{\zihao{-5} DSSAST CMS 管理员操作手册}
  \fancyhead[R]{\zihao{-5} 版本 \cmsversion}
  \fancyfoot[C]{\zihao{-5}\thepage}
  \renewcommand{\headrulewidth}{0.4pt}
}

\begin{document}
\hypersetup{pageanchor=false}
\begin{titlepage}
  \centering
  \vspace*{29mm}
  {\zihao{3}\heiti 笃实书院学生科协\par}
  \vspace{21mm}
  {\zihao{1}\bfseries DSSAST CMS\par}
  \vspace{5mm}
  {\zihao{1}\heiti 管理员操作手册\par}
  \vspace{20mm}
  {\zihao{4}适用版本\quad DSSAST CMS \cmsversion\par}
  \vspace{5mm}
  {\zihao{4}适用对象\quad 网站管理员与内容审核人员\par}
  \vfill
  {\zihao{4}二〇二六年十一月\par}
  \vspace{16mm}
\end{titlepage}

\pagenumbering{Roman}
\hypersetup{pageanchor=true}
\pagestyle{plain}
\pdfbookmark[1]{目录}{front-matter}
\begingroup
\ctexset{section={format=\centering\zihao{3}\heiti}}
\tableofcontents
\endgroup

\clearpage
\pagenumbering{arabic}
\pagestyle{manual}

\section{登录与工作台}

\subsection{登录管理后台}
在浏览器中打开 \url{https://dssast.cn/wp-login.php}，使用个人 WordPress 管理员账号登录。点击左侧的\href{https://dssast.cn/wp-admin/admin.php?page=dssast-management}{\textbf{内容管理}}打开工作台。顶部的待审核笔记数和已发布赛事数可点击，分别直达待审核队列和已发布赛事列表。

电脑上，工作台按两行两列排列，上排为\textbf{赛事管理}和\textbf{学习笔记}，下排为\textbf{新闻管理}和\textbf{公告管理}。手机上四张卡片按单列排列，可点后台顶部菜单图标打开侧栏。

\begin{manualtable}{后台侧栏与常用操作}
  \begin{tabularx}{\linewidth}{@{}p{38mm}Y@{}}
    \toprule
    \textbf{侧栏入口} & \textbf{子菜单或用途} \\
    \midrule
    内容管理 & 点击直接打开工作台。 \\
    赛事管理 & 赛事列表、添加赛事、赛事分类。 \\
    学习笔记 & 全部笔记、待审核笔记、添加笔记。 \\
    新闻管理 & 新闻列表、添加新闻、秀米同步。 \\
    公告管理 & 公告列表、添加公告。 \\
    媒体 & 查看赛事封面、新闻图片和秀米同步的媒体文件。 \\
    \bottomrule
  \end{tabularx}
\end{manualtable}

\subsection{账号与日常检查}
赛事管理、笔记审核与秀米绑定使用管理员账号。新闻和公告沿用 WordPress 原生文章权限。学生用普通账号投稿。验收时分别使用管理员、学生和访客会话。

首次接手时，在\menu{插件}{已安装插件}确认 DSSAST CMS 已启用，在\menu{设置}{常规}确认站点时区。正式站采用 \code{Asia/Shanghai}。赛事报名起止时间按站点时区填写。

\section{赛事信息管理}

\subsection{新建与发布}
\begin{enumerate}
  \item 进入\menu{赛事管理}{添加赛事}。在标题和正文中填写赛事名称、介绍、赛制、对象与参赛要求。
  \item 根据需要选择赛事分类、设置特色图片并填写摘要。赛事列表可读取摘要与封面。
  \item 在编辑页的 \textbf{赛事信息} 面板填写赛事进度、日期和报名资料，参见表~\ref{tab:competition-fields}。
  \item 保存草稿后重新打开，检查字段和正文。确认后点击 \textbf{发布}。
  \item 以访客身份打开\href{https://dssast.cn/events/}{赛事栏目}，核对赛事详情、报名时间和链接。
\end{enumerate}

\begin{manualtable}{赛事信息字段}\label{tab:competition-fields}
  \begin{tabularx}{\linewidth}{@{}p{26mm}Y@{}}
    \toprule
    \textbf{字段} & \textbf{填写要求} \\
    \midrule
    赛事状态 & 筹备中、报名中、报名结束、已结束、已取消，可选未设置。 \\
    赛事日期 & 填写一个日期，例如 2026 年 12 月 20 日。 \\
    报名开始 & 填写日期和时间，例如 2026 年 12 月 1 日 09:00。 \\
    报名截止 & 填写日期和时间，且不得早于报名开始。 \\
    报名链接 & 填写报名表的完整 HTTP 或 HTTPS 地址。 \\
    \bottomrule
  \end{tabularx}
\end{manualtable}

\subsection{维护进度和报名信息}
WordPress 的发布状态决定赛事是否公开。赛事状态表示业务进度，由管理员随报名和比赛安排手动更新。切换到“报名中”后，还应确认报名表实际开放。

所有赛事业务字段均可留空。清空字段并更新可删除原值。如果出现“赛事业务字段未保存”，整组业务字段会保留原值，此时标题或正文可能已保存。修正日期、时间顺序或链接后再次保存，重新打开编辑页逐项确认。业务字段需主动保存。

\subsection{在列表中快速编辑}
进入\menu{赛事管理}{赛事列表}，在赛事标题下点\textbf{快速编辑}。在\textbf{赛事信息}中修改赛事状态、赛事日期、报名起止时间或报名链接后，点击\textbf{更新}。遇到校验提示时，可在原面板修正后再保存。赛事介绍在完整编辑页修改，WordPress 自带的“日期”字段用于调整发布时间。

学生通过公开赛事页的报名链接进入报名表。发布后用学生账号确认表单权限。报名表负责收集报名并管理名单。停止报名时，先关闭表单提交，再更新后台赛事进度。

\section{学习笔记投稿与审核}

\subsection{学生投稿}
\begin{enumerate}
  \item 打开笔记投稿页 \url{https://dssast.cn/note-submit/}，按提示登录普通账号并返回投稿页。
  \item 填写标题和正文，点击 \textbf{提交审核}。标题输入框最多 200 个字符。正文支持基本 HTML，投稿表单暂不提供附件上传。
  \item 在 \textbf{我的投稿（最近 30 条）} 中确认稿件为待审核。新投稿每分钟最多一篇，收到频率提示时等待后再提交。
\end{enumerate}

\subsection{审核并发布}
\begin{enumerate}
  \item 进入\menu{学习笔记}{待审核笔记}，按标题或正文搜索稿件，点击 \textbf{审核}。
  \item 核对作者、正文、资料来源和外部链接。必要时编辑标题或正文。
  \item 稿件通过时，确认“退回作者修改”处于未选状态，点击 \textbf{发布}。
  \item 到\menu{学习笔记}{全部笔记}核对发布状态，并在\href{https://dssast.cn/notes/}{学习笔记栏目}或该稿件的查看链接检查公开内容。
\end{enumerate}

\par
\ifdim\dimexpr\pagegoal-\pagetotal\relax<8\baselineskip\newpage\fi
\subsection{退回、重投和再次审核}
\begin{enumerate}
  \item 在笔记的 \textbf{笔记审核} 面板填写可执行的修改意见，例如“请补充实验步骤及资料来源”。
  \item 勾选 \textbf{退回作者修改}，使用编辑页的保存按钮提交。回到全部笔记，确认该稿件已转为草稿。
  \item 作者在投稿页的 \textbf{我的投稿} 中找到该稿，点击 \textbf{修改再提交}，阅读表单上方的审核意见，修改后重新提交。
  \item 管理员在待审核队列复核。再次发布时保持退回选项未勾选。
\end{enumerate}

审核意见供管理员和稿件作者查看。作者可编辑本人已退回的草稿。请审核人员和作者主动检查队列及投稿状态。已发布笔记退回后将取消公开，操作前确认是否需要撤下。

\section{新闻公告管理}

\subsection{新建和发布}
\begin{enumerate}
  \item 发布新闻时进入\menu{新闻管理}{添加新闻}。发布公告时进入\menu{公告管理}{添加公告}。
  \item 编辑器会为文章预选对应分类，发布前请再次核对。
  \item 填写标题、正文，按需要设置特色图片和摘要。先预览，再发布或定时发布。
  \item 在对应的\textbf{新闻列表}或\textbf{公告列表}核对文章状态，再到\href{https://dssast.cn/news/}{新闻栏目}或\href{https://dssast.cn/announcements/}{公告栏目}查看发布结果。
\end{enumerate}

\subsection{分类与已有稿件}
正式站的新闻和公告分别采用 \code{news} 和 \code{announcements} 分类别名。它们用于栏目查询和归档地址。编辑已有文章时，从对应管理页搜索标题并打开编辑页，确认分类、可见性和发布时间后更新。

\section{秀米稿件同步}

\subsection{配置与绑定}
使用秀米自建网站应用向 WordPress 推送图文。技术维护人员负责在服务器配置应用 ID 和密钥，确认 HTTPS 回调可用，并为图片处理提供 GD 或 Imagick。凭据通过授权渠道交接。

截至 2026 年 11 月 23 日，DSSAST 应用在\href{https://xiumi.us/#/user/ownerpartnerbind}{秀米应用管理页}显示“无接口，不可用”。首次审核通过后，才可绑定秀米账号与同步图文。

\begin{enumerate}
  \item 打开\menu{新闻管理}{秀米同步}，将后台显示的图文、图片和通知回调地址分别填入秀米应用设置。工作台新闻卡片也有秀米同步入口。
  \item 首次审核通过后，点击 \textbf{绑定当前管理员的秀米账号}，在秀米完成授权。站点收到绑定通知后，可以用该绑定账号同步。
  \item 首次联调选用一篇测试图文。在\menu{新闻管理}{新闻列表}查找导入草稿，并在 \textbf{媒体} 检查图片。
\end{enumerate}

\begin{manualtable}{秀米回调与用途}
  \begin{tabularx}{\linewidth}{@{}p{48mm}Y@{}}
    \toprule
    \textbf{后台所示回调末段} & \textbf{用途} \\
    \midrule
    \code{/xiumi/html} & 接收图文并写入新闻草稿。 \\
    \code{/xiumi/image} & 接收图片并存入 WordPress 媒体库。 \\
    \code{/xiumi/notify} & 接收绑定、解绑、密钥重置及同步错误通知。 \\
    \bottomrule
  \end{tabularx}
\end{manualtable}

\subsection{同步、校对和发布}
在秀米选中准备同步的公众号稿件，发送到已绑定的网站。同步生成新闻草稿。管理员打开草稿，核对标题、正文、图片、封面、链接和分类，预览后手动发布。图片存入媒体库后具有公开地址，上传前应确认授权范围。

同一稿件原样重试会跳过重复写入。源文有更新时，同步可以覆盖仍处于草稿状态的正文，并保留手动调整的分类。稿件已发布或已处于其他状态时，变化的源文会被拒绝覆盖。编辑人员应在同步完毕后进行站内校对。

图片支持 JPEG、PNG、GIF 和 WebP。单张大小上限为 5 MiB，像素总数上限为 2000 万。同步后的图片经过重新编码，复杂排版与动画应在预览时复查。

\subsection{日志与验收状态}
在 \textbf{秀米同步} 页面查看最近错误，时间戳按 UTC 记录。模拟稿件已验证导入新闻草稿和重复推送去重。真实秀米联调须在应用审核通过并完成绑定后进行。具体配置见\href{https://github.com/L-B123/kx_adminPage/blob/main/dssast-cms/XIUMI.md}{秀米技术配置说明}。

\section{发布检查与常见问题}

\subsection{发布检查}
\begin{itemize}
  \item \textbf{赛事}　核对进度、日期、报名时间、链接和表单权限。
  \item \textbf{学习笔记}　核对作者与来源。退回时填写意见并确认已转草稿。
  \item \textbf{新闻与公告}　核对分类、图片、来源和发布时间，发布后查看栏目。
  \item \textbf{秀米稿件}　确认图文存为草稿，校对正文和媒体后再发布。
\end{itemize}

\begin{manualtable}{常见问题与处理}
  \begin{tabularx}{\linewidth}{@{}p{39mm}Y@{}}
    \toprule
    \textbf{现象} & \textbf{处理方法} \\
    \midrule
    后台缺少管理菜单 & 确认以管理员身份登录，检查 DSSAST CMS 是否已启用。 \\
    正文更新，赛事时间没变 & 查看业务字段保存提示，修正时间格式、起止顺序及链接后重新保存。 \\
    学生无法修改稿件 & 检查稿件状态与作者账号。需要修改的待审核稿件由管理员退回。 \\
    投稿页提示验证失败 & 复制未保存的正文，刷新页面或重新登录再提交。持续出现时请运维检查投稿页缓存。 \\
    笔记保存后仍为草稿 & 检查“退回作者修改”的勾选状态，再按审核结果操作。 \\
    新闻或公告没有出现在栏目 & 核对文章分类、发布状态及定时发布时间。 \\
    秀米认证或绑定失败 & 先确认审核通过，再查看最近错误。请技术同学核对应用凭据、HTTPS 和绑定通知。 \\
    图文已导入，图片缺失 & 核对文件格式和大小，请技术维护人员确认 PHP 图像处理扩展。 \\
    \bottomrule
  \end{tabularx}
\end{manualtable}

\section{维护与交接}

\subsection{安装与升级}
正式站安装或升级由管理员与技术维护人员协调。先在测试站从\menu{插件}{安装插件}上传 DSSAST CMS 的 ZIP 包并启用，验收菜单、内容编辑和投稿。备份正式站数据库、插件文件及上传目录后再更新。完成后到\menu{插件}{已安装插件}确认版本和启用状态，并复查发布流程。

\subsection{交接清单}
交接 WordPress 登录入口、个人账号、栏目分类、赛事更新负责人、笔记审核约定和秀米应用管理人。新管理员在测试站练习赛事发布、笔记退回与重投，并核对新闻公告分类。投稿页及其编辑链接应排除整页缓存。凭据通过授权渠道交接。

\end{document}
